% GENERATED FILE. Edit canonical lessons, then run npm run generate:books.
% canonical-source-hash: fnv1a64:10863e5868bf40d2
% canonical-lessons: BN-C260-sadh, BN-C260-puroskar, BN-C260-shasti, BN-C260-protijogita, BN-C260-joy

\chapter{Longing, Prize, Punishment, Competition, Victory}
\label{ch:bn-longing-prize-punishment-competition-victory}
\hlchaptermodality{260}

\begin{chapteropening}
\textbf{By the end of this chapter:} \emph{I can say a longing, a prize, a punishment, a competition and victory.}

Complete the last lesson of chapter 260: I can say a longing, a prize, a punishment, a competition and victory.
\end{chapteropening}

\section[\texorpdfstring{\bn{সাধ} (sādh)}{সাধ (sādh)}]{\bn{সাধ} (sādh) — a longing}

\label{lesson:BN-C260-sadh}



\begin{quote}
Before the new one: say the Bengali for an application, then the Bengali for an identity card.
\end{quote}

\subsection*{You'll want to know: \bn{সাধ}}
\textbf{\bn{সাধ}} — \emph{sādh} — \textquotedblleft{}a longing\textquotedblright{}.

\begin{cousinweb}
\textbf{\bn{সাধ}} (\emph{sādh}) is a tadbhava, a word inherited from Sanskrit and reshaped by sound change. It comes from \emph{śraddhā}, \textquotedblleft{}trust; longing\textquotedblright{} (CDIAL 12678). CDIAL lists Bengali \emph{sādh}, \textquotedblleft{}wish, fancy\textquotedblright{}.
\end{cousinweb}

\begin{grammarlens}[title={What you've built}]
One more word for this chapter.
\end{grammarlens}

\subsection*{Your turn}
\begin{itemize}
\raggedright
  \item \emph{Say it:} \emph{sādh}
  \item \emph{Say it:} \emph{sādh}, once more
  \item \emph{Say it:} say \emph{sādh}
  \item \emph{Recall:} say the Bengali for an application, then the Bengali for an identity card, then say \emph{sādh} again
\end{itemize}

\begin{tcolorbox}[breakable,colback=teal!4,colframe=teal!35!black,title={Before you move on}]
What does \bn{সাধ} mean? (\textquotedblleft{}A longing\textquotedblright{}.) Say it once more.
\end{tcolorbox}
\section[\texorpdfstring{\bn{পুরস্কার} (purôskār)}{পুরস্কার (purôskār)}]{\bn{পুরস্কার} (purôskār) — a prize}

\label{lesson:BN-C260-puroskar}



\begin{quote}
Before the new one: say the Bengali for an identity card, then the Bengali for a longing.
\end{quote}

\subsection*{You'll want to know: \bn{পুরস্কার}}
\textbf{\bn{পুরস্কার}} — \emph{purôskār} — \textquotedblleft{}a prize\textquotedblright{}.

\begin{cousinweb}
\textbf{\bn{পুরস্কার}} (\emph{purôskār}) is Sanskrit \emph{puraskāra}, \textquotedblleft{}honouring\textquotedblright{}, taken over whole as a learned loan, a tatsama.
\end{cousinweb}

\begin{grammarlens}[title={What you've built}]
One more word for this chapter.
\end{grammarlens}

\subsection*{Your turn}
\begin{itemize}
\raggedright
  \item \emph{Say it:} \emph{purôskār}
  \item \emph{Say it:} \emph{purôskār}, once more
  \item \emph{Say it:} say \emph{purôskār}
  \item \emph{Recall:} say the Bengali for an identity card, then the Bengali for a longing, then say \emph{purôskār} again
\end{itemize}

\begin{tcolorbox}[breakable,colback=teal!4,colframe=teal!35!black,title={Before you move on}]
What does \bn{পুরস্কার} mean? (\textquotedblleft{}A prize\textquotedblright{}.) Say it once more.
\end{tcolorbox}
\section[\texorpdfstring{\bn{শাস্তি} (shāsti)}{শাস্তি (shāsti)}]{\bn{শাস্তি} (shāsti) — a punishment}

\label{lesson:BN-C260-shasti}



\begin{quote}
Before the new one: say the Bengali for a longing, then the Bengali for a prize.
\end{quote}

\subsection*{You'll want to know: \bn{শাস্তি}}
\textbf{\bn{শাস্তি}} — \emph{shāsti} — \textquotedblleft{}a punishment\textquotedblright{}.

\begin{cousinweb}
\textbf{\bn{শাস্তি}} (\emph{shāsti}) is Sanskrit \emph{śāsti}, \textquotedblleft{}correction, punishment\textquotedblright{}, taken over whole as a learned loan, a tatsama.
\end{cousinweb}

\begin{grammarlens}[title={What you've built}]
One more word for this chapter.
\end{grammarlens}

\subsection*{Your turn}
\begin{itemize}
\raggedright
  \item \emph{Say it:} \emph{shāsti}
  \item \emph{Say it:} \emph{shāsti}, once more
  \item \emph{Say it:} say \emph{shāsti}
  \item \emph{Recall:} say the Bengali for a longing, then the Bengali for a prize, then say \emph{shāsti} again
\end{itemize}

\begin{tcolorbox}[breakable,colback=teal!4,colframe=teal!35!black,title={Before you move on}]
What does \bn{শাস্তি} mean? (\textquotedblleft{}A punishment\textquotedblright{}.) Say it once more.
\end{tcolorbox}
\section[\texorpdfstring{\bn{প্রতিযোগিতা} (prôtijogitā)}{প্রতিযোগিতা (prôtijogitā)}]{\bn{প্রতিযোগিতা} (prôtijogitā) — a competition}

\label{lesson:BN-C260-protijogita}



\begin{quote}
Before the new one: say the Bengali for a prize, then the Bengali for a punishment.
\end{quote}

\subsection*{You'll want to know: \bn{প্রতিযোগিতা}}
\textbf{\bn{প্রতিযোগিতা}} — \emph{prôtijogitā} — \textquotedblleft{}a competition\textquotedblright{}.

\begin{cousinweb}
\textbf{\bn{প্রতিযোগিতা}} (\emph{prôtijogitā}) is Sanskrit \emph{pratiyogitā}, \textquotedblleft{}rivalry\textquotedblright{}, taken over whole as a learned loan, a tatsama.
\end{cousinweb}

\begin{grammarlens}[title={What you've built}]
One more word for this chapter.
\end{grammarlens}

\subsection*{Your turn}
\begin{itemize}
\raggedright
  \item \emph{Say it:} \emph{prôtijogitā}
  \item \emph{Say it:} \emph{prôtijogitā}, once more
  \item \emph{Say it:} say \emph{prôtijogitā}
  \item \emph{Recall:} say the Bengali for a prize, then the Bengali for a punishment, then say \emph{prôtijogitā} again
\end{itemize}

\begin{tcolorbox}[breakable,colback=teal!4,colframe=teal!35!black,title={Before you move on}]
What does \bn{প্রতিযোগিতা} mean? (\textquotedblleft{}A competition\textquotedblright{}.) Say it once more.
\end{tcolorbox}
\section[\texorpdfstring{\bn{জয়} (jôy)}{জয় (jôy)}]{\bn{জয়} (jôy) — victory}

\label{lesson:BN-C260-joy}



\begin{quote}
Before the new one: say the Bengali for a punishment, then the Bengali for a competition.
\end{quote}

\subsection*{You'll want to know: \bn{জয়}}
\textbf{\bn{জয়}} — \emph{jôy} — \textquotedblleft{}victory\textquotedblright{}.

\begin{cousinweb}
\textbf{\bn{জয়}} (\emph{jôy}) is Sanskrit \emph{jaya}, \textquotedblleft{}victory\textquotedblright{}, taken over whole as a learned loan, a tatsama.
\end{cousinweb}

\begin{grammarlens}[title={What you've built}]
One more word for this chapter.
\end{grammarlens}

\subsection*{Your turn}
\begin{itemize}
\raggedright
  \item \emph{Say it:} \emph{jôy}
  \item \emph{Say it:} \emph{jôy}, once more
  \item \emph{Say it:} say \emph{jôy}
  \item \emph{Recall:} say the Bengali for a punishment, then the Bengali for a competition, then say \emph{jôy} again
\end{itemize}

\begin{tcolorbox}[breakable,colback=teal!4,colframe=teal!35!black,title={Before you move on}]
What does \bn{জয়} mean? (\textquotedblleft{}Victory\textquotedblright{}.) Say it once more.
\end{tcolorbox}
